\documentclass[a4paper,11pt,oneside]{book}
\usepackage[DIV=14,BCOR=2mm,headinclude=true,footinclude=false]{typearea}
\title{Sections and Chapters}
\author{Gubert Farnsworth}
\date{ }
\begin{document}

\maketitle

\tableofcontents

\newpage
\section*{Abstract}

\chapter{Introduction}


\section{Genetic Screening - its meaning and purpose}


\section{Dynamic Screening}


\section{Aim of our study}


\section{cAMP signaling pathway}


\section{Phosphodiesterases (PDEs)}


\section{Beta-adrenergic receptor, agonist and antagonist }


\subsection{Beta-adrenergic receptor}


\subsection{Beta adrenergic agonist}


\subsection{Beta adrenergic antagonist/blocker}


\section{Sensor for measuring cAMP dynamics }


\section{Förster resonance energy transfer (FRET) }


\section{FLIM (Fluorescence Lifetime Imaging)}


\section{Study model and experiment outline}


\newpage
\chapter{Materials and Methods}
\section{Cell culture}


\section{Experiments on the Fd-FLIM set up}


\section{Experiments on the confocal set up}


\section{Generating cells with PDE knockdowns}


\section{End point calibration using Forskolin }


\section{Data acquisition and analysis using ImageJ and MATLAB }


\newpage
\chapter{Results}
\section{Experiments on the Fd-FLIM}
\subsection{Testing different Isoproterenol (IP) concentrations as stimuli}


\subsection{Testing IBMX with different Isoproterenol concentrations}


\subsection{Determining the ideal concentration of Isoproterenol as a stimuli}


\section{Experiments on the Confocal}
\subsection{Determining the correct Isoproterenol concentration as stimuli}


\subsection{Determining the correct Propranolol concentration as a blocker}


\subsection{Test PDE screen by using IBMX}


\section{Screen on the Confocal}
\subsection{Performing the experiment}


\subsection{Getting the lifetime data of each cell}


\subsection{Automatically acquiring the cAMP breakdown rate per cell}


\section{Strategies for a more efficient and automated screen}
\subsection{Change in performance of the experiment}


\subsection{Automatically getting the lifetime data of each cell}


\subsection{Getting the cAMP breakdown rate per cell }


\section{Comparison of the cAMP breakdown rate for each PDE}


\section{PDE3A is the dominant PDE in HeLa cells}






\newpage
\chapter{Discussion}
\newpage

\newpage
\chapter{Supplementary material}

Hello, as said by \cite{bos2006}



\bibliographystyle{plain}
\bibliography{ReportSravasti2019_reflist}
\end{document}